% Part: turing-machines
% Chapter: undecidability
% Section: unsolvability-decision-problem

\documentclass[../../../include/open-logic-section]{subfiles}

\begin{document}

\olfileid{tur}{und}{uns}
\olsection{د پرېکړې مسئله ناحل ده}

\begin{thm}
\ollabel{thm:decision-prob}
د پرېکړې مسئله ناحل ده: داسې ټيورينګ ماشين~$D$ نشته چې پر هغې پټې
له پيلېدو وروسته چې د لومړۍ درجې منطق يو !!a{sentence}~$!B$ د ننوت
په توګه لري، $D$~بالاخره ودرېږي او هله او يوازې هله~$1$ وباسي چې~$!B$
معتبر وي، او که نۀ وي~$0$ وباسي۔
\end{thm}

\begin{proof}
فرض کړئ د پرېکړې مسئله حل کېدونکې وے، يعنې فرض کړئ چې يو ټيورينګ
ماشين~$D$ شته۔ بيا به مو د درېدنې مسئله په لاندې ډول حل کړې وے۔ يو
ټيورينګ ماشين~$E$ جوړوو چې د ټيورينګ ماشين~$M_e$ عدد~$e$ او
ننوت~$w$ ور کړل شي، اړوند !!{sentence}~$!T(M_e, w) \lif !E(M_e,
w)$ محاسبه کړي او په داسې حال کښې ودرېږي چې پر پټه تر ټولو کيڼه
مربع لولي۔ بيا به ماشين~$E \concat D$ ته چې ننوتونه~$e$ او~$w$ ور
کړل شي، لومړے~$!T(M_e, w) \lif !E(M_e, w)$ محاسبه کړي او ورپسې به
د پرېکړې مسئلې ماشين~$D$ پر همدې ننوت وچلوي۔ $D$ هله او يوازې هله
له وت~$1$ سره درېږي چې~$!T(M_e, w) \lif !E(M_e, w)$ معتبر وي، او
که نۀ وي~$0$ وباسي۔ د~\olref[ver]{lem:halt-if-valid} او
\olref[ver]{lem:valid-if-halt} له مخې، $!T(M_e, w) \lif !E(M_e, w)$
هله او يوازې هله معتبر دے چې~$M_e$ پر ننوت~$w$ ودرېږي۔ نو~$E
\concat D$ چې ننوتونه~$e$ او~$w$ ور کړل شي، هله او يوازې هله له
وت~$1$ سره درېږي چې~$M_e$ پر ننوت~$w$ ودرېږي، او که نۀ وي له
وت~$0$ سره درېږي۔ په بله وينا، $E \concat D$ به د درېدنې مسئله حل
کړې وے۔ خو د~\olref[hal]{thm:halting-problem} له مخې پوهېږو چې داسې
ټيورينګ ماشين شتون نۀ شي لرلے۔
\end{proof}

\begin{cor}\ollabel{cor:undecidable-sat}%
دا پرېکړه کېدونکې نۀ ده چې آيا د لومړۍ درجې منطق کوم خپل‌سری
!!{sentence} د صدق وړ دے۔
\end{cor}

\begin{proof}
  فرض کړئ چې د صدق وړتيا يو ټيورينګ ماشين~$S$ پرېکړه کولے شوه۔ بيا
  به مو د پرېکړې مسئله په لاندې ډول حل کړې وے: يو
  !!a{sentence}~$!B$ د ننوت په توګه واخلئ او~$!B$ يوه مربع ښي لوري
  ته يوسئ۔ بېرته مربع~$1$ ته راشئ او نښه~$\lnot$ وليکئ۔

  اوس ټيورينګ ماشين~$S$ وچلوئ۔ دا بالاخره درېږي او پر پټه يا~$1$
  وباسي (که~$\lnot !B$ د صدق وړ وي) يا~$0$ (که~$\lnot !B$ د صدق
  وړ نۀ وي)۔ که پر مربع~$1$ يوه~$\TMstroke$ وي، پاکه ئې کړئ؛ که
  مربع~$1$ تشه وي، يوه~$\TMstroke$ پرې وليکئ، او بيا ودرېږئ۔

  دا ټيورينګ ماشين تل درېږي، او وت ئې هله او يوازې هله~$1$ وي چې
  $\lnot !B$ د صدق وړ نۀ وي، او که نۀ وي~$0$ وي۔ څرنګه چې~$!B$ هله
  او يوازې هله معتبر دے چې~$\lnot !B$ د صدق وړ نۀ وي، ماشين هله او
  يوازې هله~$1$ وباسي چې~$!B$ معتبر وي، او که نۀ وي~$0$ وباسي؛
  يعنې د پرېکړې مسئله به حل کړي۔
  \textbf{د ژباړې د نسخې سمون (OLCMP-062):} سرچينې د ننوت د غونډلې
  په لومړۍ يادونه کښې د ټاکلې څيز-ژبې نښه پرېښې وه؛ دلته له ورپسې
  ټولو يادونو سره سمه بشپړه نښه کارول شوې ده۔
\end{proof}

\begin{explain}
نو داسې ټيورينګ ماشين نشته چې دې پوښتنې ته تل سم ``هو'' يا ``نه''
ځواب ورکړي: ``آيا~$!B$ د لومړۍ درجې منطق يو معتبر !!{sentence}
دے؟'' خو داسې ټيورينګ ماشين \emph{شته} چې تل سم ``هو'' ځواب
ورکوي---مګر که ځواب ``نه'' وي، هېڅ نۀ درېږي۔ دا د لومړۍ درجې منطق
د سموالي او بشپړتيا له قضيې، او له دې حقيقت څخه راځي چې
!!{derivation}s په اغېزمنه توګه شمېرل کېدے شي۔
\end{explain}

\begin{thm}
  \ollabel{thm:valid-ce}%
  د لومړۍ درجې د !!{sentence}s اعتبار نيمه د پرېکړې وړ دے: داسې
  ټيورينګ ماشين~$E$ شته چې پر هغې پټې له پيلېدو وروسته چې د لومړۍ
  درجې منطق يو !!a{sentence}~$!B$ د ننوت په توګه لري، $E$~هله او يوازې
  هله بالاخره درېږي او~$1$ وباسي چې~$!B$ معتبر وي، او که نۀ وي نۀ
  درېږي۔
\end{thm}

\begin{proof}
  د لومړۍ درجې منطق ټول ممکن !!{derivation}s د يوه اغېزمن الګوريتم
  په وسيله يو په بل پسې توليدېدے شي۔ ماشين~$E$ همدا کار کوي، او کله
  چې داسې !!a{derivation} ومومي چې~$\Proves !B$ ښيي، له وت~$1$ سره
  درېږي۔ د سموالي د قضيې له مخې، که~$E$ له وت~$1$ سره درېږي، لامل
  ئې دا دے چې~$\Entails !B$۔ د بشپړتيا د قضيې له مخې، که~$\Entails
  !B$ وي، داسې !!a{derivation} شته چې~$\Proves !B$ ښيي۔ څرنګه
  چې~$E$ ټول ممکن !!{derivation}s په منظم ډول توليدوي، بالاخره به
  داسې يو ومومي چې~$\Proves !B$ ښيي؛ نو بالاخره به له وت~$1$ سره
  ودرېږي۔
\end{proof}

\end{document}
